\chapter{Reference: full enum and key list}
\label{ch:refman-enums}

\begin{aside}
The complete set of enumerated constants \maya{} exposes.
Look here when the compiler asks for a name and you've
forgotten it.
\end{aside}

\section{SpecialKey}

\begin{itemize}[leftmargin=*]
  \item \code{Up}, \code{Down}, \code{Left}, \code{Right}.
  \item \code{Enter}, \code{Escape}, \code{Tab},
    \code{Backspace}, \code{Delete}.
  \item \code{Home}, \code{End}, \code{PageUp},
    \code{PageDown}, \code{Insert}.
  \item \code{F1} through \code{F12}.
\end{itemize}

\section{BorderStyle}

\begin{itemize}[leftmargin=*]
  \item \code{None}, \code{Round}, \code{Square},
    \code{Heavy}, \code{Double}, \code{Ascii}.
\end{itemize}

\section{MouseButton}

\begin{itemize}[leftmargin=*]
  \item \code{Left}, \code{Middle}, \code{Right}.
  \item \code{WheelUp}, \code{WheelDown},
    \code{WheelLeft}, \code{WheelRight}.
  \item \code{None} for move events with no button held.
\end{itemize}

\section{MouseKind}

\begin{itemize}[leftmargin=*]
  \item \code{Press}, \code{Release}, \code{Drag},
    \code{Move}.
\end{itemize}

\section{Cursor}

\begin{itemize}[leftmargin=*]
  \item \code{Hidden}, \code{Block}, \code{Underscore},
    \code{Bar}.
  \item \code{BlinkingBlock}, \code{BlinkingBar},
    \code{BlinkingUnderscore}.
\end{itemize}

\section{Working example}

\begin{lstlisting}[language=C++]
#include <maya/maya.hpp>

using namespace maya;
using namespace maya::dsl;

struct EnumDemo {
    struct Model { std::string last = "(none)"; };
    struct Got { std::string label; };
    struct Quit {};
    using Msg = std::variant<Got, Quit>;

    static Model init() { return {}; }
    static auto update(Model m, Msg msg)
        -> std::pair<Model, Cmd<Msg>> {
        return std::visit(overload{
            [&](Got g) { m.last = g.label; return std::pair{m, Cmd<Msg>{}}; },
            [](Quit) { return std::pair{Model{}, Cmd<Msg>::quit()}; },
        }, msg);
    }

    static Element view(const Model& m) {
        return v(
            text(std::format("last: {}", m.last)),
            blank_,
            t<"try F1..F4, arrows, page up/down"> | Dim,
            t<"q to quit"> | Dim
        ) | pad<1> | border_<Round>;
    }

    static auto subscribe(const Model&) -> Sub<Msg> {
        return key_map<Msg>({
            {'q', Quit{}},
            {SpecialKey::F1,       Got{"F1"}},
            {SpecialKey::F2,       Got{"F2"}},
            {SpecialKey::F3,       Got{"F3"}},
            {SpecialKey::F4,       Got{"F4"}},
            {SpecialKey::Up,       Got{"up"}},
            {SpecialKey::Down,     Got{"down"}},
            {SpecialKey::Left,     Got{"left"}},
            {SpecialKey::Right,    Got{"right"}},
            {SpecialKey::PageUp,   Got{"pgup"}},
            {SpecialKey::PageDown, Got{"pgdn"}},
            {SpecialKey::Home,     Got{"home"}},
            {SpecialKey::End,      Got{"end"}},
        });
    }
};

int main() { run<EnumDemo>({.title = "enums"}); }
\end{lstlisting}

\section{Notes}

\begin{itemize}[leftmargin=*]
  \item Not every key reaches every terminal; F-keys in
    particular are sometimes captured by tmux or by the
    terminal emulator itself.
  \item \code{Insert} is rarely useful; default to other
    bindings.
  \item \code{shift(...)}, \code{ctrl(...)}, \code{alt(...)}
    can modify any of these.
\end{itemize}

\section{Recap}

\begin{recap}
\begin{itemize}[leftmargin=*]
  \item One enum per concern.
  \item All keys reachable via \code{key\_map}.
  \item Modifier wrappers compose.
\end{itemize}
\end{recap}

\section*{Workshop}

\begin{enumerate}[leftmargin=*]
  \item Print a label for every key you press.
  \item Build a "vim-like" set of bindings (\code{hjkl},
    \code{gg}, \code{G}).
  \item Try every cursor style with \code{RunConfig::cursor}.
\end{enumerate}
